\documentclass[11pt]{article}
\usepackage[T1]{fontenc}
\usepackage{lmodern}
\usepackage{amsmath,amssymb,amsthm,mathtools}
\usepackage{mathrsfs}
\usepackage[margin=1in]{geometry}
\usepackage{booktabs}
\usepackage[colorlinks=true,linkcolor=blue,citecolor=blue,urlcolor=blue]{hyperref}
\usepackage{enumitem}
\allowdisplaybreaks
\setlength{\emergencystretch}{2em}
\numberwithin{equation}{section}
\newtheorem{theorem}{Theorem}[section]
\newtheorem{proposition}[theorem]{Proposition}
\newtheorem{corollary}[theorem]{Corollary}
\newtheorem{lemma}[theorem]{Lemma}
\theoremstyle{definition}
\newtheorem{definition}[theorem]{Definition}
\theoremstyle{remark}
\newtheorem{remark}[theorem]{Remark}
\DeclareMathOperator{\Rea}{Re}
\DeclareMathOperator{\Ima}{Im}
\DeclareMathOperator{\arccosh}{arccosh}
\newcommand{\RH}{\mathrm{RH}}
\newcommand{\dd}{\,d}
\newcommand{\Lag}{\mathscr L_0}
\newcommand{\Sh}{\mathcal S^\sharp}

\title{Euler-Factor Controls for the de Bruijn--Newman Constant:\protect\\
Certified and Asymptotic Lower Bounds}
\author{Drafted for Edward Baker}
\date{14 October 2026}

\begin{document}
\maketitle

\begin{abstract}
For a prime $p$ and $c>0$ we change the Riemann $\xi$-function at one
Euler factor. In the Newman normalization this replaces the Riemann kernel
$\Phi$ by $\Phi_{c,p}=\Phi+\tfrac c2[\Phi(\cdot+a)+\Phi(\cdot-a)]$ with
$a=\tfrac12\log p$, and the time-zero transform $H_0(x)$ by
$H_0(x)(1+c\cos ax)$. The associated Dirichlet series is
$\zeta(s)\bigl(1+\tfrac{2\sqrt p}{c}p^{-s}+p^{1-2s}\bigr)$. It has every
Euler factor of $\zeta$ except the one at $p$, lies in the extended Selberg
class with conductor $p^2$, and its local factor at $p$ is tempered exactly
when $c\ge1$. For $c\ge1$ the de Bruijn--Newman constant $\Lambda_{c,p}$
satisfies $\Lambda_{c,p}\le0$ if and only if the Riemann hypothesis holds.
For $0<c<1$ we prove three results.
(i)~The heat flow of $\Phi_{c,p}$ is that of $\Phi$ plus two copies of
$H_t$ at complex spectral shifts $x\pm2ita$.
(ii)~By outward-rounded interval evaluation of one Laguerre expression,
$\Lambda_{1/2,p}>0.2$ for every prime $p\le41$, and
$\Lambda_{1/2,p}>0.1787854$ for $p=43,47$.
(iii)~With $t_m=4\log(1/c)/(\log p)^2$ and
$S_2=\xi''(\tfrac12)/(2\xi(\tfrac12))$,
\[
\Lambda_{c,p}>t_m-S_2t_m^2-\tfrac1{50}\bigl(\tfrac12\log p\bigr)^{-6}
\qquad(\tfrac12\le c<1,\ p\ge29).
\]
Numerically, the lowest off-line pair of zeros collides at
$t_m-S_2t_m^2+O(t_m^3)$. Since the best refereed upper bound is
$\Lambda\le0.2$, these kernels show that any proof of an upper bound
below $\Lambda_{c,p}$ must use an input that $\Phi_{c,p}$ lacks, such as
the Euler factor of $\zeta$ at $p$ or direct information about the zeros of
$\xi$; the controls already have non-real zeros at height $\pi/a<10$. We
describe how the existing proofs supply such inputs. No new
information about the zeros of $\zeta$ is claimed.
\end{abstract}

\tableofcontents

\section{Introduction}\label{sec:intro}

\subsection{The de Bruijn--Newman constant}
Let
\begin{equation}\label{eq:Phi}
\Phi(u)=\sum_{n\ge1}\bigl(2\pi^2n^4e^{9u}-3\pi n^2e^{5u}\bigr)
\exp\bigl(-\pi n^2e^{4u}\bigr),
\end{equation}
and for an even real kernel $K$ of sufficiently rapid decay put
\begin{equation}\label{eq:Ht}
H_t[K](z)=\int_0^\infty e^{tu^2}K(u)\cos(zu)\dd u ,\qquad
H_t=H_t[\Phi].
\end{equation}
Then $H_0(z)=\tfrac18\xi(\tfrac12+\tfrac{iz}2)$, where
$\xi(s)=\tfrac12s(s-1)\pi^{-s/2}\Gamma(s/2)\zeta(s)$, and
$\partial_tH_t=-\partial_z^2H_t$. De Bruijn~\cite{debruijn} showed that
real-rootedness of $H_t$ persists as $t$ increases and that $H_t$ has only
real zeros for $t\ge\tfrac12$. Newman~\cite{newman} showed that there is a
constant $\Lambda\in(-\infty,\tfrac12]$ such that $H_t$ has only real zeros
exactly when $t\ge\Lambda$. The Riemann hypothesis
is equivalent to $\Lambda\le0$. Earlier lower bounds for $\Lambda$ came
from Lehmer pairs, pairs of unusually close zeros of
$\zeta$~\cite{csordassmithvarga}; the witnesses of Section~\ref{sec:witness}
also rest on the local behaviour of a single pair of zeros. Rodgers and
Tao~\cite{rodgerstao} proved $\Lambda\ge0$; Dobner~\cite{dobner} gave a new proof and extended it to the
extended Selberg class. On the other side, Ki, Kim and Lee~\cite{kikimlee}
proved $\Lambda<\tfrac12$, the Polymath project~\cite{polymath} proved
$\Lambda\le0.22$, and Platt and Trudgian~\cite{platttrudgian} proved
$\Lambda\le0.2$, the best bound in the refereed literature. A computer-assisted
claim $\Lambda\le0.1787854$~\cite{gomila} has not been refereed and is not
used below. The survey~\cite{newmanwu} describes the wider landscape.

\subsection{Controls}
An upper bound $\Lambda\le t_0$ needs information about $\zeta$. One way to
see which information is to exhibit kernels that share many properties
with $\Phi$ but have de Bruijn--Newman constant larger than $t_0$. Analytic
shape alone gives no upper bound: $e^{-Tu^2}\Phi(u)$ is positive, even and
rapidly decaying, and its constant is $\Lambda+T$, because
$H_t[e^{-Tu^2}K]=H_{t-T}[K]$. This paper studies kernels that are much closer to $\Phi$
arithmetically: they differ from it at one Euler factor.

Fix a prime $p$, write
\begin{equation}\label{eq:a}
a=a_p=\tfrac12\log p,
\end{equation}
and for $c>0$ define
\begin{equation}\label{eq:control}
\Phi_{c,p}(u)=\Phi(u)+\frac c2\bigl[\Phi(u+a)+\Phi(u-a)\bigr],\qquad
H^{c,p}_t=H_t[\Phi_{c,p}],
\end{equation}
and let $\Lambda_{c,p}$ be the de Bruijn--Newman constant of
$\Phi_{c,p}$ (Definition~\ref{def:Lambda}). At $t=0$,
\begin{equation}\label{eq:t0}
H^{c,p}_0(x)=H_0(x)\,(1+c\cos ax).
\end{equation}
For $0<c<1$ the multiplier lies in $[1-c,1+c]$ on the real line, so
$H^{c,p}_0$ has exactly the real zeros of $H_0$. Its additional zeros lie
on the lines $\Ima x=\pm\delta$, where
\begin{equation}\label{eq:delta}
\delta=\delta_{c,p}=\frac{\arccosh(1/c)}{a}.
\end{equation}
We also use
\begin{equation}\label{eq:tm}
\tau_0=\log\frac1c,\qquad
t_m=\frac{\tau_0}{a^2}=\frac{4\log(1/c)}{(\log p)^2},\qquad
S_2=\frac{\xi''(\tfrac12)}{2\,\xi(\tfrac12)}=0.0231049931154\ldots .
\end{equation}
The quantity $S_2$ equals $\sum_{\gamma}\gamma^{-2}$, summed over the zeros
$\tfrac12+i\gamma$ of $\xi$ with $\Rea\gamma>0$; under RH this is the sum
over the positive ordinates.

\subsection{Results}
The first theorem collects the structural facts (Section~\ref{sec:controls}
and Proposition~\ref{prop:basic}).

\begin{theorem}[structure]\label{thm:structure}
Let $p$ be prime and $c>0$.
\begin{enumerate}[label=\textup{(\alph*)}]
\item $\Phi_{c,p}$ is positive and even, and $H^{c,p}_t$ is an even entire
function of order at most one for every real $t$.
\item The Dirichlet series attached to $H^{c,p}_0$ is
$F_{c,p}(s)=\zeta(s)Q_p(p^{-s})$ with
$Q_p(X)=1+\tfrac{2\sqrt p}{c}X+pX^2$. It has positive coefficients and the
Euler factor of $\zeta$ at every prime $q\ne p$. It belongs to the extended
Selberg class with degree one and conductor $p^2$. Its local factor at $p$
is tempered (its roots lie on $|X|=p^{-1/2}$) if and only if $c\ge1$.
\item If $c\ge1$, then $\Lambda_{c,p}\ge0$, and $\Lambda_{c,p}\le0$ if and
only if RH holds.
\item If $0<c<1$, then $0<\Lambda_{c,p}\le\max(1,\delta)^2/2$, and
$\Lambda_{c,p}\le\delta^2/2$ under RH.
\end{enumerate}
\end{theorem}

The heat flow of a control is determined by the heat flow of $\Phi$.

\begin{theorem}[exact flow]\label{thm:flow}
For every real $t$ and complex $z$,
\[
H^{c,p}_t(z)=H_t(z)+\frac c2\,e^{ta^2}\Bigl[e^{-iaz}H_t(z+2ita)
+e^{iaz}H_t(z-2ita)\Bigr].
\]
\end{theorem}

To first order the bracket acts like the multiplier
$1+ce^{ta^2}\cos az$, whose extra zeros become real when
$ce^{ta^2}=1$, that is at $t=t_m$. The factor
$e^{ta^2}=\exp(\tfrac t4\log^2p)$ is the weight $b_p^t$ with which the heat
flow multiplies the $p$-th Dirichlet coefficient in the effective
approximation of~\cite{polymath}. At $t_m$ the amplified local factor
reaches the tempered boundary (Remark~\ref{rem:model}).

A single negative value of the Laguerre expression
$\Lag(f;x)=f'(x)^2-f(x)f''(x)$ proves a lower bound for
$\Lambda_{c,p}$ (Proposition~\ref{prop:witness}). Interval evaluation of
such witnesses gives the following.

\begin{theorem}[certified bounds]\label{thm:certified}
$\Lambda_{1/2,p}>0.2$ for every prime $p\le41$, and
$\Lambda_{1/2,p}>0.1787854$ for $p\in\{43,47\}$.
\end{theorem}

For large $p$ we obtain an asymptotic lower bound with explicit
constants.

\begin{theorem}[asymptotic lower bound]\label{thm:asymptotic}
Let $p$ be prime and $0<c<1$, and let $a,\tau_0,t_m,S_2$ be as above.
\begin{enumerate}[label=\textup{(\roman*)}]
\item If $\tfrac12\le c<1$ and $p\ge29$, then
$\Lambda_{c,p}>t_m-S_2t_m^2-\tfrac1{50}\,a^{-6}$.
\item If $e^{-2}\le c<1$ and $p\ge293$, then
$\Lambda_{c,p}>t_m-S_2t_m^2-\tfrac1{10}\,a^{-6}$.
\end{enumerate}
In both cases the Laguerre expression of $H^{c,p}_t$ at $x=\pi/a$ is negative
at the displayed time $t$ whenever $t\ge0$. For every fixed $c\in(0,1)$,
$\Lambda_{c,p}\ge t_m-S_2t_m^2-O(t_m^3)$ as $p\to\infty$.
\end{theorem}

For $c=\tfrac12$ the bound of Theorem~\ref{thm:asymptotic}(i) exceeds
$0.2$ for $29\le p\le37$ and exceeds $0.1787854$ for $29\le p\le47$. So the
analytic and the computational results confirm each other for
$p\in\{29,31,37,43,47\}$; for $p=41$ only the computation reaches $0.2$. Under RH, Theorem~\ref{thm:structure}(d) gives
the upper bound $\delta^2/2=t_m(1+\tfrac13(1-c)+O((1-c)^2))$, so that
\[
1-S_2t_m-\frac{1}{50\,\tau_0a^4}<\frac{\Lambda_{c,p}}{t_m}
\le1+\frac{1-c}3+O\bigl((1-c)^2\bigr),
\]
and both sides tend to $1$ if $p\to\infty$ and then $c\to1$
(Corollary~\ref{cor:sandwich}).

Numerical tracking of the zeros (Section~\ref{sec:numerics}) indicates that
the lowest off-line pair is the last of the tracked pairs to collide, at a time
$T_0(c,p)=t_m-S_2t_m^2+O(t_m^3)$. We do not prove an upper bound
$\Lambda_{c,p}\le T_0(c,p)$; that would require control of all other zeros.

\subsection{What the controls say about proofs}
Theorem~\ref{thm:certified} is a statement about specific entire functions.
It also has a methodological reading (Section~\ref{sec:screening}): any
collection of hypotheses satisfied by some $\Phi_{1/2,p}$ with $p\le41$
cannot imply $\Lambda_K\le0.2$ for all kernels $K$ satisfying them. The
shared hypotheses include positivity, evenness and the decay of the kernel;
the real zeros of $H_0$; bounded relative distortion of $H_0$ on the real
line; the Euler factors and Hecke relations at
every prime other than $p$; and membership in the extended Selberg class. The
existing proofs of upper bounds avoid this obstruction by using inputs that
the controls do not have, namely numerical verification of RH up to a large
height and effective approximations built from the Dirichlet coefficients
and small Euler factors of $\zeta$. The controls make precise that these
direct inputs, which the controls fail already at height $\pi/a\le9.1$,
cannot be replaced by structural facts shared with $\Phi_{c,p}$.

\subsection{Organization}
Section~\ref{sec:controls} gives the arithmetic of the controls.
Section~\ref{sec:Lambda} recalls the de Bruijn--Newman theory for general
kernels and proves Theorem~\ref{thm:structure}. Section~\ref{sec:flow}
proves Theorem~\ref{thm:flow}. Section~\ref{sec:witness} treats Laguerre
witnesses. Section~\ref{sec:certified} proves Theorem~\ref{thm:certified},
and Section~\ref{sec:asymptotic} proves Theorem~\ref{thm:asymptotic}.
Section~\ref{sec:numerics} reports the numerical collision times,
Section~\ref{sec:screening} discusses consequences for proof strategies, and
Section~\ref{sec:open} lists open problems.

\section{The controls}\label{sec:controls}

\begin{lemma}\label{lem:kernel}
Let $p$ be prime and $c>0$.
\begin{enumerate}[label=\textup{(\alph*)}]
\item $\Phi_{c,p}$ is positive and even on $\mathbb R$, and
$\Phi_{c,p}(u)=O\bigl(\exp(9|u|-\pi e^{4|u|-4a})\bigr)$.
\item For every real $t$, $H^{c,p}_t$ is an even entire function, real on
the real axis, of order at most one.
\item $H^{c,p}_0(x)=H_0(x)(1+c\cos ax)$.
\end{enumerate}
\end{lemma}

\begin{proof}
The function $\Phi$ is even, by the functional equation of the theta function,
and positive on $\mathbb R$ (P\'olya~\cite{polya}; see
also~\cite{csordasnorfolkvarga}). Hence $\Phi_{c,p}$ is positive, and it is even
because the bracket in~\eqref{eq:control} is. For $u\ge0$ the series
\eqref{eq:Phi} gives $0<\Phi(u)\le C e^{9u}\exp(-\pi e^{4u})$, and by
evenness this holds with $|u|$ in place of $u$. Since
$|u\pm a|\ge|u|-a$, part (a) follows.

For complex $z$ with $|\Ima z|=y$,
$|H^{c,p}_t(z)|\le\int_0^\infty e^{tu^2}\Phi_{c,p}(u)\cosh(yu)\dd u$. By
(a) the integrand is at most
$C\exp\bigl(tu^2+(9+y)u-\pi e^{4u-4a}\bigr)$, and the maximum over $u$ of
the exponent is $O(y\log y)$ as $y\to\infty$. Hence
$|H^{c,p}_t(z)|\le\exp(O(|z|\log|z|))$, which gives order at most one.
The function is even and real on $\mathbb R$ because the kernel is even and
real.

For (c), extend $\Phi$ evenly and write
$H_0[K](x)=\tfrac12\int_{\mathbb R}K(u)e^{ixu}\dd u$ for even $K$. Then
$\tfrac12\int_{\mathbb R}[\Phi(u+a)+\Phi(u-a)]e^{ixu}\dd u
=(e^{-iax}+e^{iax})H_0(x)=2\cos(ax)\,H_0(x)$.
\end{proof}

\begin{proposition}[arithmetic]\label{prop:arith}
Let $s=\tfrac12+\tfrac{ix}2$ and $Q_p(X)=1+\tfrac{2\sqrt p}{c}X+pX^2$.
\begin{enumerate}[label=\textup{(\alph*)}]
\item $1+c\cos ax=\tfrac c2\,p^{s-1/2}Q_p(p^{-s})$, so
$H^{c,p}_0(x)=\tfrac c{16}\,p^{s-1/2}\,\xi(s)\,Q_p(p^{-s})$.
\item $F_{c,p}(s)=\zeta(s)Q_p(p^{-s})=\sum_{n\ge1}a_nn^{-s}$ with
$a_n=1+\tfrac{2\sqrt p}c\mathbf 1_{p\mid n}+p\,\mathbf 1_{p^2\mid n}>0$,
and
\[
F_{c,p}(s)=\frac{Q_p(p^{-s})}{1-p^{-s}}\prod_{q\ne p}\frac1{1-q^{-s}}
\qquad(\Rea s>1).
\]
\item The completed function $\Lambda_{c,p}^{\mathrm{L}}(s)=(p/\sqrt\pi)^s\,
\Gamma(s/2)F_{c,p}(s)$ satisfies
$\Lambda^{\mathrm L}_{c,p}(s)=\Lambda^{\mathrm L}_{c,p}(1-s)$. Thus
$F_{c,p}$ belongs to the extended Selberg class $\Sh$, with degree one and
conductor $p^2$.
\item The roots $X_\pm$ of $Q_p$ satisfy $X_+X_-=1/p$. If $c\ge1$ they lie on
$|X|=p^{-1/2}$, and they coincide when $c=1$. If $0<c<1$ they are real and
negative, with $|X_\pm|=p^{-1/2}e^{\mp\arccosh(1/c)}$. Consequently the zeros
of $Q_p(p^{-s})$ lie on $\Rea s=\tfrac12$ if $c\ge1$. If $c<1$ they lie on
$\Rea s=\tfrac12\pm\arccosh(1/c)/\log p$, which in the variable $x$ are the
lines $\Ima x=\pm\delta$.
\end{enumerate}
\end{proposition}

\begin{proof}
(a) Here $p^{s-1/2}=e^{iax}$ and $p^{-s}=p^{-1/2}e^{-iax}$, so
$\tfrac c2p^{s-1/2}Q_p(p^{-s})=\tfrac c2e^{iax}+1+\tfrac c2e^{-iax}$.
Combine with~\eqref{eq:t0} and $H_0=\xi/8$.
(b) is the expansion of $\zeta(s)Q_p(p^{-s})$.
(c) By (a), $\Lambda^{\mathrm L}_{c,p}(s)
=\frac{2\sqrt p}{s(s-1)}\cdot\frac2c\,(1+c\cos ax)\,\xi(s)$. Each factor is
invariant under $s\mapsto1-s$, that is $x\mapsto-x$. The function $F_{c,p}$
is a Dirichlet series absolutely convergent for $\Rea s>1$, $(s-1)F_{c,p}$
is entire of finite order, and the functional equation has the shape required
in~\cite{dobner}. The conductor is $(2\pi)\cdot(p/\sqrt\pi)^2\cdot\tfrac12
=p^2$.
(d) The discriminant of $pX^2+\tfrac{2\sqrt p}cX+1$ is $4p(c^{-2}-1)$. If
$c\ge1$ the roots are conjugate, or equal when $c=1$, and $|X_\pm|^2=X_+X_-=1/p$.
If $c<1$, then $X_\pm=-p^{-1/2}(c^{-1}\mp\sqrt{c^{-2}-1})
=-p^{-1/2}e^{\mp\arccosh(1/c)}$. Finally $p^{-s}=X$ forces
$\Rea s=-\log|X|/\log p$, and $\Ima x=-2(\Rea s-\tfrac12)$.
\end{proof}

\begin{remark}[Hecke interpretation]\label{rem:hecke}
For the normalized Hecke operator
$T_pf(z)=p^{-1/2}\bigl(f(pz)+\sum_{j=0}^{p-1}f((z+j)/p)\bigr)$, the
Eisenstein series $E(z,\tfrac12+i\nu)$ has eigenvalue
$\lambda_p(\nu)=p^{i\nu}+p^{-i\nu}=2\cos(\nu\log p)$. With $\nu=x/2$,
\[
1+c\cos ax=1+\tfrac c2\,\lambda_p(x/2).
\]
Thus $H^{c,p}_0$ is the spectral density of $\Phi$ multiplied by the
eigenvalue of the Hecke element $1+\tfrac c2T_p$. A multiplier that depends
only on $\lambda_p$ leaves the local data at every prime $q\ne p$ unchanged;
this is the spectral form of Proposition~\ref{prop:arith}(b).
\end{remark}

\begin{remark}[the Selberg class]\label{rem:selberg}
$F_{c,p}$ has an Euler product, but its local factor at $p$,
$Q_p(p^{-s})/(1-p^{-s})$, has zeros. The coefficients of $\log F_{c,p}$ at
$p^k$ are $-(X_+^{-k}+X_-^{-k})/k$ plus the contribution of
$(1-p^{-s})^{-1}$, and they have size at least a constant times $p^{k/2}/k$
for infinitely many $k$. So $F_{c,p}$ violates the Euler-product axiom of the
Selberg class for every $c>0$, while it lies in $\Sh$, which imposes no such
axiom.
\end{remark}

\section{De Bruijn--Newman constants of the controls}\label{sec:Lambda}

\begin{lemma}\label{lem:dbn}
Let $K$ be a real even kernel with $K(u)=O(\exp(-|u|^{2+\eta}))$ for some
$\eta>0$.
\begin{enumerate}[label=\textup{(\alph*)}]
\item \textup{(de Bruijn~\cite[Theorem~13]{debruijn}; see
also~\cite[Theorem~3.2]{polymath})} If all zeros of $H_{t_0}[K]$ lie in
$|\Ima z|\le\Delta$, then for $t\ge t_0$ all zeros of $H_t[K]$ lie in
$|\Ima z|\le\sqrt{\max(\Delta^2-2(t-t_0),0)}$. In particular, real-rootedness
persists as $t$ increases.
\item The set $Z_K$ of $t\in\mathbb R$ for which $H_t[K]$ has only real zeros
is closed.
\item If $H_t[K]$ is of order less than two and has only real zeros, then
$\Lag(H_t[K];x)\ge0$ for every real $x$.
\end{enumerate}
\end{lemma}

\begin{proof}
(b) The map $t\mapsto H_t[K]$ is continuous into the entire functions with
the topology of locally uniform convergence, and $H_t[K]\not\equiv0$. By
Hurwitz's theorem, a locally uniform limit of entire functions with only real
zeros has only real zeros.
(c) By Hadamard's theorem, an entire function $f$ of order less than two,
real on $\mathbb R$, with only real zeros $x_k$, has the form
$f(z)=Cz^me^{\beta z}\prod_k(1-z/x_k)e^{z/x_k}$ with real $\beta$. For real
$x$ that is not a zero,
$-(\log f)''(x)=mx^{-2}+\sum_k(x-x_k)^{-2}\ge0$, so
$\Lag(f;x)=-f(x)^2(\log f)''(x)\ge0$. At a zero, $\Lag(f;x)=f'(x)^2\ge0$.
\end{proof}

\begin{definition}\label{def:Lambda}
For $K$ as in Lemma~\ref{lem:dbn} put $\Lambda_K=\inf Z_K\in[-\infty,\infty]$.
By Lemma~\ref{lem:dbn}, $Z_K=[\Lambda_K,\infty)$ whenever $\Lambda_K$ is
finite. We write $\Lambda=\Lambda_\Phi$ and $\Lambda_{c,p}=\Lambda_{\Phi_{c,p}}$.
\end{definition}

Lemma~\ref{lem:kernel}(a) shows that $\Phi_{c,p}$ satisfies the hypothesis of
Lemma~\ref{lem:dbn}.

\begin{proposition}\label{prop:basic}
Let $p$ be prime.
\begin{enumerate}[label=\textup{(\alph*)}]
\item If $0<c<1$, then $0<\Lambda_{c,p}\le\max(1,\delta)^2/2$, and under RH
$\Lambda_{c,p}\le\delta^2/2$.
\item If $c\ge1$, then $\Lambda_{c,p}\ge0$, and $\Lambda_{c,p}\le0$ holds if
and only if RH holds. Under RH, $\Lambda_{c,p}=0$.
\item For every $c>0$, $\Lambda_{c,p}\ge0$.
\end{enumerate}
\end{proposition}

\begin{proof}
By Lemma~\ref{lem:kernel}(c), the zeros of $H^{c,p}_0$ are those of $H_0$,
which lie in $|\Ima x|<1$ because the nontrivial zeros of $\zeta$ lie in
$0<\Rea s<1$, together with the zeros of $1+c\cos ax$.
(a) For $c<1$ the latter lie on $|\Ima x|=\delta$
(Proposition~\ref{prop:arith}(d)). Lemma~\ref{lem:dbn}(a) with
$\Delta=\max(1,\delta)$, or $\Delta=\delta$ under RH, gives the upper bounds.
Since $H^{c,p}_0$ has non-real zeros, $0\notin Z_{\Phi_{c,p}}$, hence
$\Lambda_{c,p}>0$.
(b) For $c\ge1$ the zeros of $1+c\cos ax$ are real. Hence $H^{c,p}_0$ has only
real zeros if and only if $H_0$ does, that is, if and only if RH holds. By
Definition~\ref{def:Lambda} this is equivalent to $\Lambda_{c,p}\le0$.
(c) By Proposition~\ref{prop:arith}(c), $F_{c,p}\in\Sh$. Its completed
function in the sense of~\cite[Section~2]{dobner} is a constant multiple of
$H^{c,p}_0((s-\tfrac12)\cdot\tfrac2i)$, and Dobner's heat deformation agrees
with $H^{c,p}_t$ up to the same constant. Dobner's theorem~\cite[Theorem~2]{dobner}
gives $\Lambda_{c,p}\ge0$. Together with (b) this proves the remaining claims.
\end{proof}

Theorem~\ref{thm:structure} is Lemma~\ref{lem:kernel},
Proposition~\ref{prop:arith} and Proposition~\ref{prop:basic}.

\section{The exact flow}\label{sec:flow}

\begin{proof}[Proof of Theorem~\ref{thm:flow}]
Extend $\Phi$ evenly and write
$N_t[K](z)=\tfrac12\int_{\mathbb R}e^{tu^2}K(u)e^{izu}\dd u$, so that
$N_t[\Phi]=H_t$ and $H^{c,p}_t=H_t+\tfrac c2(N_t[\Phi(\cdot+a)]
+N_t[\Phi(\cdot-a)])$. Substituting $v=u\pm a$,
\[
N_t[\Phi(\cdot\pm a)](z)
=\tfrac12\int_{\mathbb R}e^{t(v\mp a)^2}\Phi(v)e^{iz(v\mp a)}\dd v
=e^{ta^2}e^{\mp iaz}\,\tfrac12\int_{\mathbb R}e^{tv^2}\Phi(v)
e^{i(z\pm2ita)v}\dd v ,
\]
which is $e^{ta^2}e^{\mp iaz}H_t(z\pm2ita)$. All integrals converge
absolutely by Lemma~\ref{lem:kernel}(a).
\end{proof}

For real $x$, $H_t(x-2ita)=\overline{H_t(x+2ita)}$, so
\begin{equation}\label{eq:flowreal}
H^{c,p}_t(x)=H_t(x)+C\,\Rea\bigl[e^{-iax}H_t(x+2ita)\bigr],\qquad
C=ce^{ta^2}.
\end{equation}

\begin{remark}[the model multiplier]\label{rem:model}
Replacing $H_t(x\pm2ita)$ by $H_t(x)$ in Theorem~\ref{thm:flow} gives
$H_t(x)(1+C\cos ax)$. With $C=ce^{ta^2}$, the extra zeros of the multiplier
are non-real for $C<1$, double for $C=1$ and real for $C>1$. They become real
at $t=t_m=\tau_0/a^2$. By Proposition~\ref{prop:arith}(d), the condition
$C\ge1$ says that the local factor $Q_p$ with $c$ replaced by $C$ is
tempered. The amplification $e^{ta^2}=\exp(\tfrac t4\log^2p)$ is the weight
$b_p^t=\exp(\tfrac t4\log^2p)$ that the heat flow attaches to the $p$-th
Dirichlet coefficient in the effective approximation of~\cite{polymath}. To
first order, the control collides when that weight has amplified its local
perturbation to the tempered boundary. The shifts $\pm2ita$ produce the
correction $-S_2t_m^2$ of Theorem~\ref{thm:asymptotic}.
\end{remark}

\section{Laguerre witnesses}\label{sec:witness}

\begin{proposition}[witness criterion]\label{prop:witness}
If $\Lag(H^{c,p}_t;x_0)<0$ for some real $x_0$, then $H^{c,p}_t$ has a
non-real zero and $\Lambda_{c,p}>t$.
\end{proposition}

\begin{proof}
If $t\ge\Lambda_{c,p}$, then $t\in Z_{\Phi_{c,p}}$ by
Definition~\ref{def:Lambda}, so $H^{c,p}_t$ has only real zeros. It has
order at most one by Lemma~\ref{lem:kernel}(b), so
Lemma~\ref{lem:dbn}(c) gives $\Lag(H^{c,p}_t;x)\ge0$ for all real $x$.
\end{proof}

Near an isolated pair of non-real zeros the Laguerre expression has a simple
local form.

\begin{proposition}\label{prop:local}
Suppose that $H(x)=((x-X)^2+y^2)\,g(x)$ near a real point $X$, with $g$
smooth and real and $y>0$. Then
\[
\Lag(H;X)=-2y^2g(X)^2+y^4\Lag(g;X),\qquad
\frac{\Lag(H;X)}{H(X)^2}=-\frac2{y^2}+\frac{\Lag(g;X)}{g(X)^2}.
\]
\end{proposition}

\begin{proof}
With $q(x)=(x-X)^2+y^2$ one has $q(X)=y^2$, $q'(X)=0$ and $q''=2$. Then
$H(X)=y^2g$, $H'(X)=y^2g'$ and $H''(X)=2g+y^2g''$, evaluated at $X$. So
$H'^2-HH''=y^4(g'^2-gg'')-2y^2g^2$.
\end{proof}

\begin{remark}\label{rem:margin}
As a pair approaches an ordinary collision, with $g$, $g'$, $g''$ and $1/g$
locally bounded, a witness exists at the real part of the pair. The
normalized value tends to $-\infty$ like $-2/y^2$, but the raw value
$\Lag(H;X)=-2y^2g(X)^2+O(y^4)$ tends to $0$. A certification must therefore
resolve a raw negative number, not merely a large normalized ratio. In the
computations of Section~\ref{sec:certified} the raw values are at least
$5\cdot10^{-6}$ in absolute value, and their enclosures are narrower than
$10^{-51}$.

Far from a collision, witnesses need not exist, since the other zeros
contribute positive terms $(x-x_k)^{-2}$ to $\Lag/H^2$. For example,
for $(p,c)=(2,0.05)$ at $t=0$ the off-line zero lies at about
$9.065+10.642i$. A numerical scan nevertheless found $\Lag(H^{c,p}_0)>0$
at $1200$ points of $(0,60]$.
\end{remark}

\section{Certified lower bounds}\label{sec:certified}

\subsection{Statement}
Theorem~\ref{thm:certified} follows from Proposition~\ref{prop:witness} and
the following computation.

\begin{proposition}\label{prop:table}
For $c=\tfrac12$ and each row of Table~\ref{tab:certified}, the enclosure of
$\Lag(H^{1/2,p}_t;x_0)$ is contained in $(-\infty,0)$. The table lists the
upper endpoint, rounded upward.
\end{proposition}

\begin{table}[ht]
\centering
\begin{tabular}{rllrrr}
\toprule
$p$ & $t$ & $x_0$ & $H^{1/2,p}_t(x_0)$ & upper bound for $\Lag$ & $\Lag/H^2$ \\
\midrule
2 & $0.2$ & 9.0194 & $1.8797\cdot10^{-2}$ & $-4.046\cdot10^{-5}$ & $-0.1145$ \\
3 & $0.2$ & 5.6906 & $2.4165\cdot10^{-2}$ & $-1.945\cdot10^{-4}$ & $-0.3332$ \\
5 & $0.2$ & 3.8844 & $2.4561\cdot10^{-2}$ & $-5.144\cdot10^{-4}$ & $-0.8529$ \\
7 & $0.2$ & 3.2128 & $2.3184\cdot10^{-2}$ & $-7.795\cdot10^{-4}$ & $-1.450$ \\
11 & $0.2$ & 2.6072 & $1.9911\cdot10^{-2}$ & $-1.149\cdot10^{-3}$ & $-2.900$ \\
13 & $0.2$ & 2.4374 & $1.8303\cdot10^{-2}$ & $-1.268\cdot10^{-3}$ & $-3.786$ \\
17 & $0.2$ & 2.2066 & $1.5239\cdot10^{-2}$ & $-1.395\cdot10^{-3}$ & $-6.011$ \\
19 & $0.2$ & 2.1232 & $1.3781\cdot10^{-2}$ & $-1.411\cdot10^{-3}$ & $-7.433$ \\
23 & $0.2$ & 1.9939 & $1.0997\cdot10^{-2}$ & $-1.359\cdot10^{-3}$ & $-11.24$ \\
29 & $0.2$ & 1.8566 & $7.0966\cdot10^{-3}$ & $-1.094\cdot10^{-3}$ & $-21.74$ \\
31 & $0.2$ & 1.8206 & $5.8577\cdot10^{-3}$ & $-9.624\cdot10^{-4}$ & $-28.05$ \\
37 & $0.2$ & 1.7314 & $2.2927\cdot10^{-3}$ & $-4.443\cdot10^{-4}$ & $-84.53$ \\
41 & $0.2$ & 1.6835 & $2.4913\cdot10^{-5}$ & $-5.291\cdot10^{-6}$ & $-8525$ \\
43 & $0.1787854$ & 1.6622 & $3.4657\cdot10^{-3}$ & $-7.142\cdot10^{-4}$ & $-59.47$ \\
47 & $0.1787854$ & 1.6238 & $1.6833\cdot10^{-3}$ & $-3.750\cdot10^{-4}$ & $-132.4$ \\
\bottomrule
\end{tabular}
\caption{Certified Laguerre witnesses for $c=\tfrac12$. The values of $H$
and $\Lag/H^2$ are rounded midpoints of their enclosures. The enclosures
of $H$ are narrower than $2\cdot10^{-50}$, those of $\Lag$ narrower than
$10^{-51}$, and those of $\Lag/H^2$ narrower than $10^{-41}$. Each $x_0$ is
a point near the real part of the lowest off-line pair.}
\label{tab:certified}
\end{table}

For $p=41$ the pair is close to its collision. Tracking gives
$T_0(\tfrac12,41)\approx0.2001167$ (Section~\ref{sec:numerics}). The
normalized witness is large, but the raw value is only $-5.3\cdot10^{-6}$, as
<<<<<<< HEAD
Remark~\ref{rem:margin} predicts. A separate certificate of this case, with
different arithmetic and quadrature, is described in the acknowledgements;
at its point $x=1.68413$ the method of this section gives enclosures inside
those of the separate certificate. For $p\ge43$ the lowest pair has already
=======
Remark~\ref{rem:margin} predicts. For $p\ge43$ the lowest pair has already
>>>>>>> 94745f3 (Merged work from ChatGPT)
collided by $t=0.2$, so the bound is certified at $t=0.1787854$ instead.

\subsection{Method}
All integrands below are even, so
$\int_0^\infty F=\tfrac12\int_{\mathbb R}F$. We use the trapezoid rule on the
real line.

\begin{lemma}[Trefethen--Weideman~\cite{trefethenweideman}, Theorem~5.1]
\label{lem:trap}
Let $F$ be analytic in the strip $|\Ima u|<b$, with $F(u)\to0$ uniformly as
$|u|\to\infty$ in the strip, and $\int_{\mathbb R}|F(u+iv)|\dd u\le M$ for
$|v|<b$. Then for every $h>0$,
\[
\Bigl|\int_{\mathbb R}F(u)\dd u-h\sum_{k\in\mathbb Z}F(kh)\Bigr|
\le\frac{2M}{e^{2\pi b/h}-1}.
\]
\end{lemma}

\begin{lemma}[strip majorant]\label{lem:majorant}
The function $\Phi$ extends analytically to the strip $|\Ima w|<\pi/8$, where
it is even. For $0\le b<\pi/8$, $\kappa=\cos4b$ and $|\Ima w|\le b$,
\[
|\Phi(w)|\le\widetilde\Phi_\kappa(|\Rea w|),\qquad
\widetilde\Phi_\kappa(r)=\sum_{n\ge1}\bigl(2\pi^2n^4e^{9r}+3\pi n^2e^{5r}\bigr)
\exp\bigl(-\pi\kappa n^2e^{4r}\bigr).
\]
\end{lemma}

\begin{proof}
For $|\Ima w|<\pi/8$ we have $\Rea e^{4w}=e^{4\Rea w}\cos(4\Ima w)>0$, and
\[
\bigl|\exp(-\pi n^2e^{4w})\bigr|=\exp\bigl(-\pi n^2e^{4\Rea w}\cos(4\Ima w)\bigr).
\]
Hence the series~\eqref{eq:Phi} converges locally uniformly on the strip and
defines an analytic function there, which equals $\Phi$ on $\mathbb R$ and is
therefore even. For $\Rea w\ge0$ the displayed identity gives the bound
termwise, using $\cos(4\Ima w)\ge\cos4b=\kappa$. For $\Rea w<0$ apply the bound
to $-w$.
\end{proof}

\begin{lemma}[tails]\label{lem:tails}
Let $\kappa\in[\tfrac13,1]$ and $r\ge2$. Then
$\widetilde\Phi_\kappa(r)\le2(2\pi^2+3\pi)\exp(9r-\pi\kappa e^{4r})$, and
the right side decreases in $r$. For $N\ge8$ and
$r\ge0$,
\[
\sum_{n>N}\bigl(2\pi^2n^4e^{9r}+3\pi n^2e^{5r}\bigr)e^{-\pi\kappa n^2e^{4r}}
\le\frac{(2\pi^2+3\pi)9^4e^{9r}e^{-81\pi\kappa e^{4r}}}
{1-(10/9)^4e^{-19\pi\kappa e^{4r}}} .
\]
\end{lemma}

\begin{proof}
Use $3\pi n^2e^{5r}\le3\pi n^4e^{9r}$ for $r\ge0$. The ratio of consecutive
terms of $n^4e^{-\pi\kappa n^2E}$ is at most $(10/9)^4e^{-19\pi\kappa E}$ for
$n\ge9$. For the first bound, the factor
$\sum_n n^4e^{-\pi\kappa(n^2-1)e^{4r}}$ is at most $2$ when $r\ge2$, because
$\pi\kappa e^{8}>3000$; and $9-4\pi\kappa e^{4r}<0$ there.
\end{proof}

\medskip\noindent\emph{The computation.} Fix $c,p,t,x_0$. The functions
\[
F_0(u)=e^{tu^2}\Phi_{c,p}(u)\cos x_0u,\quad
F_1(u)=-ue^{tu^2}\Phi_{c,p}(u)\sin x_0u,\quad
F_2(u)=-u^2e^{tu^2}\Phi_{c,p}(u)\cos x_0u
\]
are even, and $\int_0^\infty F_j$ equals $H$, $H'$ and $H''$ at $x_0$.
Take $b=3/10$, so that $\kappa=\cos(6/5)>0.36$, and $h=1/64$. For
$|\Ima u|\le b$, Lemma~\ref{lem:majorant} gives
\[
|F_j(u+iv)|\le\cosh(x_0b)\,(u^2+b^2)^{j/2}e^{tu^2}
\Bigl[\widetilde\Phi_\kappa(|u|)+\tfrac c2\bigl(\widetilde\Phi_\kappa(|u+a|)
+\widetilde\Phi_\kappa(|u-a|)\bigr)\Bigr].
\]
We bound $M_j$ by an interval Riemann upper sum of this majorant over
$[0,U]$, with $U\ge a+2$ and cells of width $1/16$, doubled by evenness. The
contribution of $|u|>U$ is bounded with Lemma~\ref{lem:tails}, using
$e^{4s}\ge1+4s+8s^2$; it is below $10^{-1400}$. The discretization error of
Lemma~\ref{lem:trap} for the half-line integral is then at most
$M_j/(e^{2\pi b/h}-1)$, which is below $10^{-50}$ in every case. The
trapezoid sum is truncated at $U$, and the omitted terms are bounded by the
same tail estimate. On the grid $u=k/64$ the shifted thetas use the exact
relations $e^{4(u\pm a)}=p^{\pm2}e^{4u}$, $e^{5(u\pm a)}=p^{\pm5/2}e^{5u}$
and $e^{9(u\pm a)}=p^{\pm9/2}e^{9u}$. For $u<a$, the evenness of $\Phi$ is
used. Eight terms of each theta series are summed explicitly, and the rest is
bounded by Lemma~\ref{lem:tails}.

\medskip\noindent\emph{Arithmetic.} Every quantity is an interval with $60$-digit
decimal endpoints, rounded outward at each operation. Python's
\texttt{decimal} module computes $\exp$ and $\log$ correctly rounded, and
we widen the result by one unit in the last place. Negation is exact. Any
operation performed outside the two directed-rounding contexts must be
exact: the default context is set to raise an exception on any inexact or
rounded result. Square roots of integers are verified by squaring. We
compute $\pi$ by Machin's formula with alternating-series remainders, and
$\sin$ and $\cos$ by Taylor series after reduction by $2\pi$, with Lagrange
remainders. As consistency checks, every case was recomputed with $h=1/32$,
and the two enclosures of each of $H,H',H''$ overlap; the enclosure of $\pi$
is compared with $70$ known decimals; and the enclosure of
$m_0(0)=\xi(\tfrac12)/8$ of Lemma~\ref{lem:moments} is compared with an
independent value of $\xi(\tfrac12)$. The program uses only the Python
standard library and runs in about $45$ seconds; its location is given in
the acknowledgements.

\section{An asymptotic lower bound}\label{sec:asymptotic}

\subsection{Moments}
For $t\ge0$ and $r\ge0$ put
\begin{equation}\label{eq:moments}
m_k(t)=\int_0^\infty e^{tu^2}\Phi(u)u^k\dd u,\qquad
\mu_4(t,r)=\int_0^\infty e^{tu^2}\Phi(u)u^4\cosh(ru)\dd u,
\end{equation}
and $k_t=m_2(t)/m_0(t)$. Since $\Phi>0$, all of these are positive and
increase in $t$ and $r$. Moreover $H_t(0)=m_0(t)$ and $H_t''(0)=-m_2(t)$.
Since $H_0(z)=\xi(\tfrac12+\tfrac{iz}2)/8$, we get $k_0=\xi''(\tfrac12)/(4\xi(\tfrac12))$, that
is,
\begin{equation}\label{eq:S2k0}
S_2=2k_0 .
\end{equation}
Also $\tfrac{d}{dt}k_t=(m_4m_0-m_2^2)/m_0^2\in[0,m_4/m_0]$ by the
Cauchy--Schwarz inequality.

\begin{lemma}\label{lem:moments}
The following enclosures hold.
\begin{center}
\begin{tabular}{lll}
\toprule
quantity & value & enclosure width \\
\midrule
$m_0(0)=\xi(\tfrac12)/8$ & $0.06214009727353926373\ldots$ & $<2\cdot10^{-51}$\\
$m_2(0)$ & $7.17873259848294922976\ldots\cdot10^{-4}$ & $<10^{-51}$\\
$S_2=2m_2(0)/m_0(0)$ & $0.02310499311541897078\ldots$ & $<10^{-50}$\\
$m_2(\tfrac14)$ & $7.23696857055556974027\ldots\cdot10^{-4}$ & $<10^{-51}$\\
$m_4(\tfrac14)$ & $2.34423514799060488665\ldots\cdot10^{-5}$ & $<10^{-52}$\\
$\mu_4(\tfrac14,3)$ & $2.90765558946641021959\ldots\cdot10^{-5}$ & $<10^{-51}$\\
\bottomrule
\end{tabular}
\end{center}
Consequently, for $0\le t\le\tfrac14$ and $0\le r\le3$,
\begin{equation}\label{eq:constants}
k_t\le\bar k:=0.0116463,\qquad
k_t-k_0\le\bar d\,t,\quad \bar d:=0.000378,\qquad
\frac{\mu_4(t,r)}{m_0(t)}\le\bar\mu:=0.00046792 .
\end{equation}
\end{lemma}

\begin{proof}
The enclosures are computed by the method of Section~\ref{sec:certified}
with $\Phi$ in place of $\Phi_{c,p}$. The weights $u^k$ and $\cosh(ru)$ are
bounded on the strip by $(u^2+b^2)^{k/2}$ and $e^{r|u|}$. The constants
follow from monotonicity: $m_0(t)\ge m_0(0)$, $m_2(t)\le m_2(\tfrac14)$,
$m_4(t)\le m_4(\tfrac14)$ and $\mu_4(t,r)\le\mu_4(\tfrac14,3)$. The program
checks that the rounded values in~\eqref{eq:constants} are upper bounds for
the certified ratios, and it uses these rounded values below.
\end{proof}

\subsection{Local expansion}
Let $A(z)=H_t(z)/H_t(0)$. Expanding the cosine,
\begin{equation}\label{eq:Aexp}
A(z)=1-\tfrac12k_tz^2+\rho(z),\qquad
\rho(z)=\frac1{m_0(t)}\int_0^\infty e^{tu^2}\Phi(u)
\bigl[\cos zu-1+\tfrac12(zu)^2\bigr]\dd u .
\end{equation}

\begin{lemma}\label{lem:remainder}
For complex $z$ with $|z|\le R$,
\[
|\rho(z)|\le\frac{|z|^4}{24}\,\mu,\qquad
|\rho'(z)|\le\frac{|z|^3}{6}\,\mu,\qquad
|\rho''(z)|\le\frac{|z|^2}{2}\,\mu,\qquad
\mu=\frac{\mu_4(t,R)}{m_0(t)} .
\]
\end{lemma}

\begin{proof}
With $y=|zu|$: $|\cos zu-1+\tfrac12(zu)^2|\le\cosh y-1-\tfrac12y^2\le
\tfrac1{24}y^4\cosh y$. Next, $|{-u}\sin zu+zu^2|\le|u|(\sinh y-y)\le
\tfrac16|u|y^3\cosh y$. Finally, $|u^2(1-\cos zu)|\le u^2(\cosh y-1)\le
\tfrac12u^2y^2\cosh y$. Each comparison is termwise in the power series.
Integrate against the positive weight $e^{tu^2}\Phi(u)$ and use
$\cosh y\le\cosh(Ru)$.
\end{proof}

\subsection{The witness at \texorpdfstring{$x_0=\pi/a$}{x0 = pi/a}}
Fix $0<c<1$, a prime $p$, and $0\le t<t_m$. Put $\tau=ta^2$, $C=ce^\tau
=e^{\tau-\tau_0}<1$ and $\varepsilon=a^{-2}$. For real $\theta$ set
$x=(\pi+\theta)/a$, so that $e^{-iax}=-e^{-i\theta}$. By~\eqref{eq:flowreal},
\begin{equation}\label{eq:F}
F(\theta):=\frac{H^{c,p}_t\bigl((\pi+\theta)/a\bigr)}{H_t(0)}
=A\Bigl(\frac{\pi+\theta}a\Bigr)
-C\,\Rea\Bigl[e^{-i\theta}A\Bigl(\frac{\pi+\theta+2i\tau}a\Bigr)\Bigr].
\end{equation}
Since $\Lag(H^{c,p}_t;\pi/a)=a^2m_0(t)^2\bigl(F'(0)^2-F(0)F''(0)\bigr)$,
a witness at $\pi/a$ means $F(0)F''(0)>F'(0)^2$. Write $x_0=\pi/a$ and
$w_0=(\pi+2i\tau)/a$. Differentiating~\eqref{eq:F},
\begin{align*}
F(0)&=A(x_0)-C\Rea A(w_0),\\
F'(0)&=A'(x_0)/a-C\bigl[\Ima A(w_0)+\Rea A'(w_0)/a\bigr],\\
F''(0)&=A''(x_0)/a^2+C\Rea A(w_0)-2C\Ima A'(w_0)/a-C\Rea A''(w_0)/a^2 .
\end{align*}
Inserting~\eqref{eq:Aexp}, with $k=k_t$,
\begin{align}
F(0)&=(1-C)\bigl(1-\tfrac12k\pi^2\varepsilon\bigr)-2Ck\tau^2\varepsilon
+\rho(x_0)-C\Rea\rho(w_0),\label{eq:F0}\\
F'(0)&=-k\pi\varepsilon(1-C-2C\tau)+\rho'(x_0)/a-C\Ima\rho(w_0)
-C\Rea\rho'(w_0)/a,\label{eq:F1}\\
F''(0)&=C-k\varepsilon(1-C)-\tfrac12Ck\pi^2\varepsilon
+2Ck\tau^2\varepsilon+4Ck\tau\varepsilon+R_2,\label{eq:F2}
\end{align}
where
$R_2=\rho''(x_0)\varepsilon+C\Rea\rho(w_0)-2C\Ima\rho'(w_0)/a
-C\Rea\rho''(w_0)\varepsilon$.

\begin{lemma}\label{lem:Fbounds}
Suppose $\tau_0\le\bar\tau_0$, $\varepsilon\le\bar\varepsilon$,
$\bar\tau_0\bar\varepsilon\le\tfrac14$ and
$(\pi^2+4\bar\tau_0^2)\bar\varepsilon\le9$. Put $P=\pi^2+4\tau_0^2$ and
$\sigma=\tau_0-\tau$. Then, with the constants of~\eqref{eq:constants},
\begin{align*}
F(0)&\ge\bigl(\sigma-\tfrac12\sigma^2\bigr)\bigl(1-\tfrac12\bar k\pi^2\varepsilon\bigr)
-2k_0\tau_0^2\varepsilon-2\bar d\tau_0^3\varepsilon^2
-\tfrac1{24}\bar\mu(\pi^4+P^2)\varepsilon^2,\\
|F'(0)|&\le\varepsilon D_1,\qquad D_1=\bar k\pi(1+2\tau_0)
+\bar\mu\bigl(\tfrac{\pi^3}6+\tfrac{P^2}{24}+\tfrac{P^{3/2}}6\bigr)\varepsilon,\\
F''(0)&\ge e^{-\sigma}-\bar k\varepsilon\bigl(1+\tfrac{\pi^2}2\bigr)
-\bar\mu D_2\varepsilon^2,\qquad
D_2=\tfrac{\pi^2}2+\tfrac{P^2}{24}+\tfrac{P^{3/2}}3+\tfrac P2 .
\end{align*}
\end{lemma}

\begin{proof}
Here $t=\tau\varepsilon\le\tau_0\varepsilon\le\tfrac14$, and
$|x_0|\le|w_0|=\sqrt{(\pi^2+4\tau^2)\varepsilon}\le\sqrt{P\varepsilon}\le3$.
So Lemma~\ref{lem:remainder} applies with $R=3$ and
$\mu\le\bar\mu$. Also $x_0^4=\pi^4\varepsilon^2$,
$|w_0|^4\le P^2\varepsilon^2$, $x_0^3/a=\pi^3\varepsilon^2$,
$|w_0|^3/a\le P^{3/2}\varepsilon^2$, $x_0^2\varepsilon=\pi^2\varepsilon^2$
and $|w_0|^2\varepsilon\le P\varepsilon^2$. In~\eqref{eq:F0} use
$1-C=1-e^{-\sigma}\ge\sigma-\tfrac12\sigma^2$, $k\le\bar k$, $C\le1$,
$\tau\le\tau_0$ and $k\le k_0+\bar dt\le k_0+\bar d\tau_0\varepsilon$.
In~\eqref{eq:F1} use $|1-C-2C\tau|\le1+2\tau_0$. In~\eqref{eq:F2} drop the
two nonnegative terms $2Ck\tau^2\varepsilon$ and $4Ck\tau\varepsilon$, and
use $(1-C)+\tfrac12C\pi^2\le1+\tfrac12\pi^2$.
\end{proof}

\subsection{Proof of Theorem~\ref{thm:asymptotic}}
Let $Q>0$ and choose $\sigma=\varepsilon s$ with
$s=2k_0\tau_0^2+Q\varepsilon$, that is,
\[
t=\frac{\tau_0-\sigma}{a^2}=t_m-2k_0t_m^2-Q\varepsilon^3
=t_m-S_2t_m^2-Q\,a^{-6},
\]
by~\eqref{eq:S2k0}. If $t<0$ there is nothing to prove, since
$\Lambda_{c,p}>0$. Otherwise Lemma~\ref{lem:Fbounds} gives, after dropping
the nonnegative term $\tfrac14\bar k\pi^2\sigma^2\varepsilon$,
\[
\frac{F(0)}{\varepsilon^2}\ge Q-\frac{s^2}2-\frac{\bar k\pi^2}2s
-2\bar d\tau_0^3-\frac{\bar\mu}{24}(\pi^4+P^2) .
\]
Write $F_2=e^{-\varepsilon s}-\bar k\varepsilon(1+\tfrac{\pi^2}2)
-\bar\mu D_2\varepsilon^2$. Then $F(0)F''(0)>F'(0)^2$ follows from $F_2>0$
and
\begin{equation}\label{eq:W}
Q>\frac{s^2}2+\frac{\bar k\pi^2}2s+2\bar d\tau_0^3
+\frac{\bar\mu}{24}(\pi^4+P^2)+\frac{D_1^2}{F_2}.
\end{equation}
Indeed, \eqref{eq:W} gives $F(0)>\varepsilon^2D_1^2/F_2\ge F'(0)^2/F''(0)$.
Every term on the right of~\eqref{eq:W} increases with $\tau_0$ and with
$\varepsilon$, and $F_2$ decreases with both. So~\eqref{eq:W} holds for all
$\tau_0\le\bar\tau_0$ and $\varepsilon\le\bar\varepsilon$ once it holds at
$(\bar\tau_0,\bar\varepsilon)$, with $k_0$ replaced by an upper bound.

For (i) take $\bar\tau_0=\log2$ and $\bar\varepsilon=4/\log^229$, so
$\bar\tau_0\bar\varepsilon<0.2446$ and $(\pi^2+4\bar\tau_0^2)\bar\varepsilon
<4.16$. With $Q=\tfrac1{50}$, interval arithmetic gives $F_2>0.967$ and the
right side of~\eqref{eq:W} below $0.01449$.
For (ii) take $\bar\tau_0=2$ and $\bar\varepsilon=4/\log^2293$, so
$\bar\tau_0\bar\varepsilon<0.2480$ and $(\pi^2+16)\bar\varepsilon<3.21$.
With $Q=\tfrac1{10}$, interval arithmetic gives $F_2>0.977$ and the right
side below $0.0680$. In both cases $\Lag(H^{c,p}_t;\pi/a)<0$, and
Proposition~\ref{prop:witness} gives $\Lambda_{c,p}>t$.

For fixed $c$ and $p\to\infty$, $\varepsilon\to0$. The hypotheses of
Lemma~\ref{lem:Fbounds} then hold with $\bar\tau_0=\tau_0$, and the right
side of~\eqref{eq:W} tends to a finite limit. Any larger $Q$ works for all
large $p$, which gives the $O(t_m^3)$ statement, since
$a^{-6}=t_m^3/\tau_0^3$. \qed

\begin{corollary}[under RH]\label{cor:sandwich}
Assume RH. For $\tfrac12\le c<1$ and $p\ge29$,
\[
t_m-S_2t_m^2-\tfrac1{50}a^{-6}<\Lambda_{c,p}\le\frac{\delta^2}2
=t_m\,g(c),\qquad g(c)=\frac{\arccosh^2(1/c)}{2\log(1/c)}
=1+\frac{1-c}3+O\bigl((1-c)^2\bigr).
\]
\end{corollary}

\begin{proof}
Combine Theorem~\ref{thm:asymptotic}(i) with
Proposition~\ref{prop:basic}(a), and note $\delta^2/2=
\arccosh^2(1/c)/(2a^2)$. The expansion of $g$ follows from
$\arccosh(1+\eta)^2=2\eta-\tfrac13\eta^2+O(\eta^3)$ and
$\log(1+\eta)=\eta-\tfrac12\eta^2+O(\eta^3)$ with $1/c=1+\eta$.
\end{proof}

\section{Numerical collision times}\label{sec:numerics}

This section is numerical evidence, not a certificate. We followed the
$k$-th off-line zero of $H^{c,p}_t$ from its position $(2k+1)\pi/a+i\delta$
at $t=0$ for $k=0,1,2$, with adaptive steps in $t$. Near the axis we solved
$H=H'=0$ by Newton's method in $(t,x)$. The integrals were computed with
the trapezoid rule on the contour $\Ima u=\tfrac14$, which was stable to
about thirteen digits under grid refinement, and Theorem~\ref{thm:flow}
was checked against direct quadrature to relative accuracy $10^{-13}$. Let
$T_0(c,p)$ denote the computed collision time of the lowest pair ($k=0$).

\begin{table}[ht]
\centering\small
\begin{tabular}{rrrrrrr}
\toprule
$p$ & $c$ & $t_m$ & $T_0(c,p)$ & $(t_m-T_0)/(S_2t_m^2)$ & Thm.~\ref{thm:asymptotic} & $\delta^2/2$\\
\midrule
2 & 0.5 & 5.7708 & 5.04322 & 0.946 & -- & 7.220\\
3 & 0.5 & 2.2972 & 2.17782 & 0.979 & -- & 2.874\\
5 & 0.5 & 1.0704 & 1.04416 & 0.990 & -- & 1.339\\
11 & 0.5 & 0.48220 & 0.476849 & 0.9957 & -- & 0.6033\\
29 & 0.5 & 0.244525 & 0.243146 & 0.998 & 0.242265 & 0.3059\\
41 & 0.5 & 0.201049 & 0.200117 & 0.9982 & 0.199627 & 0.2515\\
43 & 0.5 & 0.195989 & 0.195103 & 0.9983 & 0.194650 & 0.2452\\
47 & 0.5 & 0.187038 & 0.186231 & 0.998 & 0.185837 & 0.2340\\
101 & 0.5 & 0.130173 & 0.129781 & 0.99884 & 0.129649 & 0.1629\\
1009 & 0.5 & 0.057954 & 0.0578767 & 0.99948 & 0.0578650 & 0.0725\\
10007 & 0.5 & 0.032679 & 0.0326543 & 0.99971 & 0.0326522 & 0.0409\\
2 & 0.9 & 0.87718 & 0.857882 & 1.085 & -- & 0.908\\
101 & 0.9 & 0.019787 & 0.0197776 & 1.0017 & 0.0196451 & 0.0205\\
10007 & 0.9 & 0.0049673 & 0.00496673 & 1.0004 & 0.0049646 & 0.00514\\
\bottomrule
\end{tabular}
\caption{Computed lowest-pair collision times (floating point) and the
rigorous lower bound of Theorem~\ref{thm:asymptotic}, where it applies. The
last column is the upper bound for $\Lambda_{c,p}$ under RH.}
\label{tab:numerics}
\end{table}

The computations show the following.
\begin{enumerate}[label=(\arabic*)]
\item In the twelve cases in which three pairs were tracked (all rows
except $p=29$ and $p=47$), the lowest pair collides last among them. The higher pairs do not collide in order of $k$: for
$p=3$ the pair near $x=28.6$, next to the first zero of $\xi$, collides
first. Pairs with $k\ge3$ were not tracked.
\item For $c=\tfrac12$ we also computed $T_0(\tfrac12,p)$ for every prime
$p\le47$; Table~\ref{tab:numerics} lists some of these values. They
decrease in $p$ and exceed $0.2$ exactly for $p\le41$, with
$T_0(\tfrac12,43)\approx0.19510$ and $T_0(\tfrac12,47)\approx0.18623$.
\item The ratio $(t_m-T_0)/(S_2t_m^2)$ tends to $1$, consistent with
$T_0=t_m-S_2t_m^2+O(t_m^3)$. The next correction has opposite signs for
$c=\tfrac12$ and $c=0.9$.
\item The collision height $x^*$ approaches $2\pi/\log p=\pi/a$; for
$(p,c)=(10007,\tfrac12)$, $x^*=0.6816$ against $\pi/a=0.68214$.
\end{enumerate}

\begin{remark}[formal upper half]
The derivation of Section~\ref{sec:asymptotic} also predicts the time at which
the lowest pair actually collides. By~\eqref{eq:F0}--\eqref{eq:F2}, the two
equations $F(\theta)=F'(\theta)=0$ have, formally, a solution with
$\theta=O(\varepsilon)$ and $1-C=2k_0\tau_0^2\varepsilon+O(\varepsilon^2)$.
Since $1-C=1-e^{\tau-\tau_0}$, this is
$T_0=t_m-S_2t_m^2+O(t_m^3)$, in agreement with Table~\ref{tab:numerics}.
Turning this into an upper bound for $\Lambda_{c,p}$ would require showing
that every other zero of $H^{c,p}_t$, including those inherited from $H_0$,
is real by time $T_0$. We do not attempt this.
\end{remark}

\section{What upper-bound proofs must use}\label{sec:screening}

\begin{corollary}\label{cor:screen}
Let $\mathcal P$ be a collection of properties of real even kernels, and
suppose every property in $\mathcal P$ holds for $\Phi_{1/2,p}$ for some prime
$p\le41$. Then $\mathcal P$ does not imply $\Lambda_K\le0.2$. If the
properties hold for $\Phi_{1/2,43}$ or $\Phi_{1/2,47}$, then $\mathcal P$
does not imply $\Lambda_K\le0.1787854$.
\end{corollary}

\begin{proof}
Otherwise Theorem~\ref{thm:certified} would be contradicted with
$K=\Phi_{1/2,p}$.
\end{proof}

The content of the corollary is the list of properties that the controls
share with $\Phi$. For $0<c<1$ and every prime $p$, Sections~\ref{sec:controls}
and~\ref{sec:Lambda} show:
\begin{enumerate}[label=(S\arabic*)]
\item $\Phi_{c,p}$ is positive and even, and decays like $\exp(-\pi e^{4|u|-4a})$.
\item $H^{c,p}_0$ has exactly the real zeros of $H_0$, with the same
multiplicities, and $1-c\le H^{c,p}_0/H_0\le1+c$ on $\mathbb R$.
\item The Dirichlet coefficients are positive and multiplicative, and the
Euler factor at every prime $q\ne p$ is that of $\zeta$; equivalently,
every Hecke relation at $q\ne p$ holds (Remark~\ref{rem:hecke}).
\item The Dirichlet series lies in $\Sh$ with the gamma factor $\Gamma(s/2)$
of $\zeta$ and conductor $p^2$, and $\Lambda_{c,p}\ge0$, as for $\zeta$.
\item The heat flow is determined by the heat flow of $\Phi$
(Theorem~\ref{thm:flow}).
\end{enumerate}
They do not share:
\begin{enumerate}[label=(D\arabic*)]
\item the Euler factor of $\zeta$ at $p$, equivalently conductor one, or the
full level-one modularity of the theta function;
\item real-rootedness of $H_0$: $H^{c,p}_0$ has the non-real zeros
$\pm(2k+1)\pi/a\pm i\delta$ for every $k\ge0$, so at every height; the
lowest have real part $\pi/a=2\pi/\log p\le9.1$;
\item the values of $\xi$ away from the zeros of $\cos ax$.
\end{enumerate}

The existing proofs of $\Lambda\le0.22$ and $\Lambda\le0.2$
\cite{polymath,platttrudgian} combine three ingredients. The first is a
numerical verification of RH up to a large height, which is an input of type
(D2) and (D3). The second is a barrier computation by direct evaluation of the
genuine heat-flowed function, again (D3). The third is an effective
approximation at large height built from the Dirichlet coefficients of
$\zeta$, with Euler mollifiers at small primes, which is (D1). The controls
fail the first input already at height $\pi/a\le9.1$, where their lowest
non-real pair lies. So these proofs are consistent with
Corollary~\ref{cor:screen}. The
corollary says something about arguments that avoid such inputs.

A structural argument for $\Lambda\le0.2$ uses no direct evaluation of $\xi$.
Examples are identities, positivity statements or symmetries of the theta
function and its lifts. Such an argument must use, for each prime $p\le41$, an
input that distinguishes $\zeta$ from $\Phi_{1/2,p}$, such as the Euler factor
of $\zeta$ at $p$ or level-one modularity. Finitely many Euler factors or Hecke
relations cannot suffice for $\Lambda\le0$ on their own: for every finite set
$S$ of primes and every $p\notin S$, the controls $\Phi_{c,p}$ with $c<1$ share
the factors at $S$ and have $\Lambda_{c,p}>0$. Hypotheses beyond those shared
by the controls, for instance quantitative information about $\xi$, are not
affected by this observation.

\begin{remark}
For $c\ge1$ the controls are RH-equivalent twins of $\Phi$
(Theorem~\ref{thm:structure}(c)). Under RH, $\Lambda_{c,p}=0$ for $c\ge1$
and $\Lambda_{c,p}>0$ for $0<c<1$, so the transition between
$\Lambda_{c,p}>0$ and $\Lambda_{c,p}\le0$ in the family is governed by the
temperedness of one local factor. This is the sense in which the screen of
Corollary~\ref{cor:screen} is arithmetic and not analytic.
\end{remark}

\section{Open problems}\label{sec:open}
\begin{enumerate}[label=(\arabic*)]
\item Prove $\Lambda_{c,p}\le t_m-S_2t_m^2+O(t_m^3)$, at least under RH, by
showing that all other zeros of $H^{c,p}_t$ are real by the time the
lowest pair collides.
\item The witnesses of Theorem~\ref{thm:certified} lie at height below $10$,
near the lowest off-line pair, which numerically collides last. Construct
controls whose last collision occurs at large height, for example from
Hecke polynomials at several primes, so as to screen arguments that use
only information at large height.
\item Extend the construction to Dirichlet $L$-functions and to other members
of $\Sh$, where Dobner's theorem gives $\Lambda\ge0$ throughout.
\item Determine the minimal arithmetic input. Does any upper bound below
$\max_{p\le41}\Lambda_{1/2,p}$ follow from the Euler factors at the primes
up to $41$ together with properties (S1)--(S5)?
\end{enumerate}

\section*{Acknowledgements and research provenance}\label{sec:provenance}
\addcontentsline{toc}{section}{Acknowledgements and research provenance}
This manuscript was drafted for Edward Baker with substantial assistance
from Claude Opus 5.5 (Anthropic; configured identifier
\texttt{claude-opus-5-5}), working in Claude Cowork at reasoning effort
max. The assistant contributed the mathematical derivations, the computer
programs, the drafting and the internal checks. The Euler-factor controls
were introduced in an internal review of the investigation portfolio in the
repository \texttt{ebbaker/shifted-zeta-positivity}, and their first study is
recorded in \texttt{papers/investigations-2/active/19\_euler\_factor\_controls}.
An internal review there by GPT-6 Astra (Codex) corrected two points of that
<<<<<<< HEAD
first study, and those corrections are used here. GPT-6 Astra also certified
$\Lambda_{1/2,41}>0.2$ independently, at $x=168413/100000$, with Arb ball
arithmetic (through python-flint) and validated quadrature of a rearranged
integral; this is Note~2 of the same folder. Evaluated at that point, the
program of Section~\ref{sec:certified} gives enclosures of $H$, $H'$, $H''$
and $\Lag$, of widths below $10^{-50}$, that lie inside the enclosures of
Note~2, whose widths are below $10^{-27}$. The two programs share no code,
arithmetic library or quadrature rule. Tracked collision times
=======
first study, and those corrections are used here. Tracked collision times
>>>>>>> 94745f3 (Merged work from ChatGPT)
are numerical estimates of the last collision of a branch, not values of
$\Lambda_{c,p}$. Near an isolated collision the raw witness margin shrinks,
as in Remark~\ref{rem:margin}.

The certificates of Theorem~\ref{thm:certified} and Lemma~\ref{lem:moments},
and the inequality~\eqref{eq:W} in Theorem~\ref{thm:asymptotic}, are produced
by a standard-library program in the \texttt{numerics} subfolder of that
folder, which also writes a JSON record. The floating-point data of
Section~\ref{sec:numerics} come from a second program in the same
subfolder:
\begin{center}\small
\begin{tabular}{l}
\texttt{check\_euler\_control\_witnesses\_claude\_20261014.py}\\
\texttt{EULER\_CONTROL\_WITNESS\_CERTIFICATES\_20261014.json}\\
\texttt{check\_euler\_factor\_controls\_claude\_20261013.py}
\end{tabular}
\end{center}
An internal review by a separate Claude instance recomputed every witness
and moment with independent high-precision quadrature. It found that the
first version of the certificate program rounded negation in Python's
default $28$-digit context, so that its enclosures were not rigorous,
although the computed values were correct to at least $24$ significant
digits. The program was
corrected as described in Section~\ref{sec:certified} and the certificates
were regenerated; the new enclosures contain the independent $75$-digit
values. These checks, and the parallel internal reviews, are not
independent specialist validation. Independent review of the arguments and
the code is still needed. No claim of literature priority is made.

\begingroup\raggedright
\begin{thebibliography}{99}

\bibitem{csordasnorfolkvarga}
G.~Csordas, T.~S.~Norfolk and R.~S.~Varga,
\emph{The Riemann hypothesis and the Tur\'an inequalities},
Trans. Amer. Math. Soc. \textbf{296} (1986), 521--541.

\bibitem{csordassmithvarga}
G.~Csordas, W.~Smith and R.~S.~Varga,
\emph{Lehmer pairs of zeros, the de Bruijn--Newman constant $\Lambda$, and
the Riemann hypothesis},
Constr. Approx. \textbf{10} (1994), 107--129.

\bibitem{debruijn}
N.~G.~de Bruijn,
\emph{The roots of trigonometric integrals},
Duke Math. J. \textbf{17} (1950), 197--226.

\bibitem{dobner}
A.~Dobner,
\emph{A proof of Newman's conjecture for the extended Selberg class},
Acta Arith. \textbf{201} (2021), 29--62. arXiv:2005.05142.

\bibitem{gomila}
J.~Gomila,
\emph{A computer-assisted upper bound $\Lambda\le0.1787854$ for the
de Bruijn--Newman constant}, announcement and audit repository
\texttt{judegomila/dbn-lambda-01787854-candidate-audit}, 2026; not peer
reviewed.
\url{https://www.judegomila.com/posts/riemann-lambda-0.1787854}.

\bibitem{kikimlee}
H.~Ki, Y.-O.~Kim and J.~Lee,
\emph{On the de Bruijn--Newman constant},
Adv. Math. \textbf{222} (2009), 281--306.

\bibitem{newman}
C.~M.~Newman,
\emph{Fourier transforms with only real zeros},
Proc. Amer. Math. Soc. \textbf{61} (1976), 245--251.

\bibitem{newmanwu}
C.~M.~Newman and W.~Wu,
\emph{Constants of de Bruijn--Newman type in analytic number theory and
statistical physics},
Bull. Amer. Math. Soc. \textbf{57} (2020), 595--614.

\bibitem{platttrudgian}
D.~Platt and T.~Trudgian,
\emph{The Riemann hypothesis is true up to $3\cdot10^{12}$},
Bull. Lond. Math. Soc. \textbf{53} (2021), 792--797.

\bibitem{polya}
G.~P\'olya,
\emph{Bemerkung \"uber die Integraldarstellung der Riemannschen
$\xi$-Funktion},
Acta Math. \textbf{48} (1926), 305--317.

\bibitem{polymath}
D.~H.~J.~Polymath,
\emph{Effective approximation of heat flow evolution of the Riemann $\xi$
function, and a new upper bound for the de Bruijn--Newman constant},
Res. Math. Sci. \textbf{6} (2019), Paper No.~31.

\bibitem{rodgerstao}
B.~Rodgers and T.~Tao,
\emph{The de Bruijn--Newman constant is non-negative},
Forum Math. Pi \textbf{8} (2020), e6.

\bibitem{trefethenweideman}
L.~N.~Trefethen and J.~A.~C.~Weideman,
\emph{The exponentially convergent trapezoidal rule},
SIAM Rev. \textbf{56} (2014), 385--458.

\end{thebibliography}
\endgroup
\end{document}
